% GENERATED FILE. Edit canonical lessons, then run npm run generate:books.
% canonical-source-hash: fnv1a64:d771a91785b21708
% canonical-lessons: SA-C407-yanta, SA-C407-valga, SA-C407-kasa, SA-C407-pratodah, SA-C407-yunakti

\chapter{Drivers and Reins}
\label{ch:sa-drivers-and-reins}
\hlchaptermodality{407}

\begin{chapteropening}
\textbf{By the end of this chapter:} \emph{I can say how a driver harnesses and handles horses.}

Complete the last lesson of chapter 407: I can say how a driver harnesses and handles horses.
\end{chapteropening}

\section[\texorpdfstring{\sk{यन्ता} (yantā)}{यन्ता (yantā)}]{\sk{यन्ता} (yantā) — a driver of horses; a charioteer}

\label{lesson:SA-C407-yanta}



\begin{quote}
Before the new one: say the Sanskrit for a queue, then the Sanskrit for free space.
\end{quote}

\subsection*{You'll want to know: \sk{यन्ता}}
\textbf{\sk{यन्ता}} — \emph{yantā} — \textquotedblleft{}a driver of horses; a charioteer\textquotedblright{}.

\begin{quote}
\sk{यन्ता} \sk{रथं} \sk{मार्गे} \sk{चालयति।}
\end{quote}

\begin{quote}
\emph{yantā rathaṃ mārge cālayati.}
\end{quote}

\begin{quote}
The charioteer drives the chariot along the road.
\end{quote}

A masculine noun.

\begin{cousinweb}
\textbf{\sk{यन्ता}} (\emph{yantā}) is built on the root √yam, \textquotedblleft{}to hold, to restrain\textquotedblright{}, with the agent ending \emph{-tṛ}; MW gives \emph{yantṛ} as a driver of horses or elephants, a charioteer (MW).
\end{cousinweb}

\begin{grammarlens}[title={What you've built}]
One more word for getting around.
\end{grammarlens}

\subsection*{Your turn}
\begin{itemize}
\raggedright
  \item \emph{Say it:} \emph{yantā}
  \item \emph{Say it:} \emph{yantā}, once more
  \item \emph{Say it:} \emph{yantā rathaṃ mārge cālayati.} — The charioteer drives the chariot along the road.
  \item \emph{Recall:} say the Sanskrit for a queue, then the Sanskrit for free space, then say \emph{yantā} again
\end{itemize}

\begin{tcolorbox}[breakable,colback=teal!4,colframe=teal!35!black,title={Before you move on}]
What does \sk{यन्ता} mean? (\textquotedblleft{}A driver of horses; a charioteer\textquotedblright{}.) Say it once more.
\end{tcolorbox}
\section[\texorpdfstring{\sk{वल्गा} (valgā)}{वल्गा (valgā)}]{\sk{वल्गा} (valgā) — a rein}

\label{lesson:SA-C407-valga}



\begin{quote}
Before the new one: say the Sanskrit for free space, then the Sanskrit for a driver of horses.
\end{quote}

\subsection*{You'll want to know: \sk{वल्गा}}
\textbf{\sk{वल्गा}} — \emph{valgā} — \textquotedblleft{}a rein\textquotedblright{}.

\begin{quote}
\sk{यन्ता} \sk{हस्ते} \sk{वल्गां} \sk{धारयति।}
\end{quote}

\begin{quote}
\emph{yantā haste valgāṃ dhārayati.}
\end{quote}

\begin{quote}
The charioteer holds the rein in his hand.
\end{quote}

A feminine noun.

\begin{cousinweb}
\textbf{\sk{वल्गा}} (\emph{valgā}) is the word MW gives for a bridle or rein; MW gives no derivation, so its origin is uncertain (MW).
\end{cousinweb}

\begin{grammarlens}[title={What you've built}]
One more word for getting around.
\end{grammarlens}

\subsection*{Your turn}
\begin{itemize}
\raggedright
  \item \emph{Say it:} \emph{valgā}
  \item \emph{Say it:} \emph{valgā}, once more
  \item \emph{Say it:} \emph{yantā haste valgāṃ dhārayati.} — The charioteer holds the rein in his hand.
  \item \emph{Recall:} say the Sanskrit for free space, then the Sanskrit for a driver of horses, then say \emph{valgā} again
\end{itemize}

\begin{tcolorbox}[breakable,colback=teal!4,colframe=teal!35!black,title={Before you move on}]
What does \sk{वल्गा} mean? (\textquotedblleft{}A rein\textquotedblright{}.) Say it once more.
\end{tcolorbox}
\section[\texorpdfstring{\sk{कशा} (kaśā)}{कशा (kaśā)}]{\sk{कशा} (kaśā) — a whip}

\label{lesson:SA-C407-kasa}



\begin{quote}
Before the new one: say the Sanskrit for a driver of horses, then the Sanskrit for a rein.
\end{quote}

\subsection*{You'll want to know: \sk{कशा}}
\textbf{\sk{कशा}} — \emph{kaśā} — \textquotedblleft{}a whip\textquotedblright{}.

\begin{quote}
\sk{यन्ता} \sk{कशया} \sk{अश्वं} \sk{न} \sk{ताडयति।}
\end{quote}

\begin{quote}
\emph{yantā kaśayā aśvaṃ na tāḍayati.}
\end{quote}

\begin{quote}
The driver does not strike the horse with the whip.
\end{quote}

A feminine noun.

\begin{cousinweb}
\textbf{\sk{कशा}} (\emph{kaśā}) is the word MW gives for a whip, from the Rigveda on; MW connects it with a root √kaś but gives no meaning for it, so its origin is uncertain (MW).
\end{cousinweb}

\begin{grammarlens}[title={What you've built}]
One more word for getting around.
\end{grammarlens}

\subsection*{Your turn}
\begin{itemize}
\raggedright
  \item \emph{Say it:} \emph{kaśā}
  \item \emph{Say it:} \emph{kaśā}, once more
  \item \emph{Say it:} \emph{yantā kaśayā aśvaṃ na tāḍayati.} — The driver does not strike the horse with the whip.
  \item \emph{Recall:} say the Sanskrit for a driver of horses, then the Sanskrit for a rein, then say \emph{kaśā} again
\end{itemize}

\begin{tcolorbox}[breakable,colback=teal!4,colframe=teal!35!black,title={Before you move on}]
What does \sk{कशा} mean? (\textquotedblleft{}A whip\textquotedblright{}.) Say it once more.
\end{tcolorbox}
\section[\texorpdfstring{\sk{प्रतोदः} (pratodaḥ)}{प्रतोदः (pratodaḥ)}]{\sk{प्रतोदः} (pratodaḥ) — a goad; a long whip}

\label{lesson:SA-C407-pratodah}



\begin{quote}
Before the new one: say the Sanskrit for a rein, then the Sanskrit for a whip.
\end{quote}

\subsection*{You'll want to know: \sk{प्रतोदः}}
\textbf{\sk{प्रतोदः}} — \emph{pratodaḥ} — \textquotedblleft{}a goad; a long whip\textquotedblright{}.

\begin{quote}
\sk{कृषकः} \sk{प्रतोदेन} \sk{वृषभं} \sk{नुदति।}
\end{quote}

\begin{quote}
\emph{kṛṣakaḥ pratodena vṛṣabhaṃ nudati.}
\end{quote}

\begin{quote}
The farmer drives the bull on with a goad.
\end{quote}

A masculine noun.

\begin{cousinweb}
\textbf{\sk{प्रतोदः}} (\emph{pratodaḥ}) is built from \emph{pra}, \textquotedblleft{}forth\textquotedblright{}, and the root √tud, \textquotedblleft{}to strike, to push\textquotedblright{}; MW gives \emph{pratoda} as a goad or long whip (MW).
\end{cousinweb}

\begin{grammarlens}[title={What you've built}]
One more word for getting around.
\end{grammarlens}

\subsection*{Your turn}
\begin{itemize}
\raggedright
  \item \emph{Say it:} \emph{pratodaḥ}
  \item \emph{Say it:} \emph{pratodaḥ}, once more
  \item \emph{Say it:} \emph{kṛṣakaḥ pratodena vṛṣabhaṃ nudati.} — The farmer drives the bull on with a goad.
  \item \emph{Recall:} say the Sanskrit for a rein, then the Sanskrit for a whip, then say \emph{pratodaḥ} again
\end{itemize}

\begin{tcolorbox}[breakable,colback=teal!4,colframe=teal!35!black,title={Before you move on}]
What does \sk{प्रतोदः} mean? (\textquotedblleft{}A goad; a long whip\textquotedblright{}.) Say it once more.
\end{tcolorbox}
\section[\texorpdfstring{\sk{युनक्ति} (yunakti)}{युनक्ति (yunakti)}]{\sk{युनक्ति} (yunakti) — harnesses; yokes}

\label{lesson:SA-C407-yunakti}



\begin{quote}
Before the new one: say the Sanskrit for a whip, then the Sanskrit for a goad.
\end{quote}

\subsection*{You'll want to know: \sk{युनक्ति}}
\textbf{\sk{युनक्ति}} — \emph{yunakti} — \textquotedblleft{}harnesses; yokes\textquotedblright{}.

\begin{quote}
\sk{यन्ता} \sk{रथे} \sk{अश्वं} \sk{युनक्ति।}
\end{quote}

\begin{quote}
\emph{yantā rathe aśvaṃ yunakti.}
\end{quote}

\begin{quote}
The charioteer harnesses the horse to the chariot.
\end{quote}

Past (he …): \sk{सः} \sk{रथे} \sk{अश्वं} \sk{युक्तवान्}, \emph{saḥ rathe aśvaṃ yuktavān}.

\begin{cousinweb}
\textbf{\sk{युनक्ति}} (\emph{yunakti}) comes from the root √yuj, \textquotedblleft{}to yoke, to join\textquotedblright{}, which also gives \emph{yojayati}, \textquotedblleft{}joins\textquotedblright{}; MW gives its participle \emph{yukta} as yoked or harnessed; MW compares Greek \emph{zeugnumi} and Latin \emph{jungo} (MW).
\end{cousinweb}

\begin{grammarlens}[title={What you've built}]
One more word for getting around.
\end{grammarlens}

\subsection*{Your turn}
\begin{itemize}
\raggedright
  \item \emph{Say it:} \emph{yunakti}
  \item \emph{Say it:} \emph{yunakti}, once more
  \item \emph{Say it:} \emph{yantā rathe aśvaṃ yunakti.} — The charioteer harnesses the horse to the chariot.
  \item \emph{Recall:} say the Sanskrit for a whip, then the Sanskrit for a goad, then say \emph{yunakti} again
\end{itemize}

\begin{tcolorbox}[breakable,colback=teal!4,colframe=teal!35!black,title={Before you move on}]
What does \sk{युनक्ति} mean? (\textquotedblleft{}Harnesses; yokes\textquotedblright{}.) Say it once more.
\end{tcolorbox}
